\documentclass[11pt,a4paper]{article}
\usepackage[utf8]{inputenc}
\usepackage[T1]{fontenc}
\usepackage[spanish]{babel}
\usepackage{graphicx}
\usepackage{float}
\usepackage[margin=2.2cm]{geometry}
\usepackage{enumitem}
\usepackage[hidelinks]{hyperref}

\title{Guía de observación\\EDA 04: parcelas y rendimiento observado}
\author{Proyecto AgroCebada}
\date{}

\newcommand{\edafigure}[3]{
  \subsection{#1}
  \begin{figure}[H]
    \centering
    \includegraphics[width=0.96\textwidth,height=0.66\textheight,keepaspectratio]{\detokenize{#2}}
    \caption{#1}
  \end{figure}
  \textbf{Preguntas guía}
  #3
}

\begin{document}
\maketitle
\tableofcontents
\newpage

\section{Fuentes y propósito}
La ubicación y forma de las parcelas provienen de
\texttt{data/original/02\_conjunto\_70\_30/Parcelas\_Reto\_AGC\_CONJUNTO.shp}.
El shapefile usa el sistema EPSG:4326 y guarda el identificador como
\texttt{ID\_POLIGON}, nombre truncado por la limitación del formato. El script lo
renombra a \texttt{ID\_POLIGONO} y lo une uno a uno con
\texttt{ID\_area\_rendimiento\_70\_30\_Reto\_AgroCebada.csv}.

El rendimiento corresponde al ciclo abril--octubre de 2025. Está observado para
138 parcelas de entrenamiento y oculto para 59 parcelas de predicción. El gris
en los mapas significa rendimiento oculto, no rendimiento cero.

\section{Exploración visual}

\edafigure{Rendimiento observado por parcela}
{../../../outputs/eda/01_tablas_maestras/04_parcelas_rendimiento/01_mapa_rendimiento_parcelas.png}{
\begin{itemize}[nosep]
  \item ¿Los rendimientos altos y bajos forman patrones espaciales?
  \item ¿Las parcelas grises están distribuidas entre las mismas regiones que las observadas?
  \item ¿Puede la ubicación explicar parte de la variación del rendimiento?
\end{itemize}}

\edafigure{Disponibilidad espacial del rendimiento}
{../../../outputs/eda/01_tablas_maestras/04_parcelas_rendimiento/02_mapa_disponibilidad_rendimiento.png}{
\begin{itemize}[nosep]
  \item ¿El conjunto de predicción ocupa regiones sin representación en entrenamiento?
  \item ¿Existe concentración espacial de etiquetas faltantes?
  \item ¿La proporción 70/30 se conserva aproximadamente por estado?
\end{itemize}}

\edafigure{Distribución global del rendimiento}
{../../../outputs/eda/01_tablas_maestras/04_parcelas_rendimiento/03_distribucion_rendimiento.png}{
\begin{itemize}[nosep]
  \item ¿Cuál es el centro y rango de los rendimientos observados?
  \item ¿La distribución es simétrica, multimodal o presenta valores extremos?
  \item ¿Hay evidencia para transformar la variable objetivo?
\end{itemize}}

\edafigure{Rendimiento por estado}
{../../../outputs/eda/01_tablas_maestras/04_parcelas_rendimiento/04_rendimiento_por_estado.png}{
\begin{itemize}[nosep]
  \item ¿Las diferencias entre estados son mayores que la variación interna?
  \item ¿Qué estado tiene mayor rendimiento y cuál menor?
  \item ¿Cómo afecta esto a la estrategia de validación?
\end{itemize}}

\edafigure{Rendimiento medio por municipio}
{../../../outputs/eda/01_tablas_maestras/04_parcelas_rendimiento/05_rendimiento_por_municipio.png}{
\begin{itemize}[nosep]
  \item ¿Qué municipios concentran rendimientos altos y bajos?
  \item ¿Cuáles promedios se apoyan en muy pocas parcelas?
  \item ¿Municipio debe emplearse como predictor, grupo de validación o ambos?
\end{itemize}}

\edafigure{Área frente a rendimiento}
{../../../outputs/eda/01_tablas_maestras/04_parcelas_rendimiento/06_area_vs_rendimiento.png}{
\begin{itemize}[nosep]
  \item ¿Las parcelas grandes tienen sistemáticamente mayor rendimiento por hectárea?
  \item ¿La aparente relación cambia por estado?
  \item ¿El área puede estar confundida con la ubicación?
\end{itemize}}

\section{Respuestas y hallazgos}

\subsection{Mapa de rendimiento}
Los valores altos se concentran principalmente en Puebla y en algunos municipios
de Hidalgo, mientras que Tlaxcala presenta valores bajos y compactos. Esto muestra
una estructura espacial clara: ubicación, clima, altitud, suelo o manejo pueden
explicar una fracción importante del objetivo. Las parcelas grises aparecen en
los tres estados y dentro de regiones que también contienen entrenamiento; no
forman una región completamente nueva.

La separación observada/oculta conserva aproximadamente el diseño 70/30 dentro
de cada estado: Hidalgo tiene 36 observadas y 15 ocultas; Puebla, 71 y 29;
Tlaxcala, 31 y 15. Por tanto, existe representación geográfica general, aunque
esto no garantiza cobertura suficiente para cada municipio pequeño.

\subsection{Distribución del objetivo}
Entre las 138 parcelas observadas, el rendimiento medio es 3.96 t/ha y la mediana
4.34 t/ha. El rango es 2.00--5.59, con cuartiles 3.20 y 4.55 y desviación estándar
0.85. La media menor que la mediana refleja una cola hacia rendimientos bajos y
la combinación de poblaciones regionales. No hay evidencia suficiente para exigir
una transformación logarítmica; primero debe evaluarse el objetivo en su escala
original, que además facilita interpretar MAE y RMSE en t/ha.

\subsection{Estado y municipio}
Puebla tiene el rendimiento medio más alto, 4.48 t/ha, seguido de Hidalgo con
3.76 y Tlaxcala con 2.98. Las medianas son 4.52, 3.50 y 3.00. Las diferencias
entre estados son grandes respecto a la dispersión interna, especialmente en
Tlaxcala, cuya desviación es 0.38.

Por municipio destacan Cuautepec de Hinojosa (4.60), Almoloya (4.52) y
Chignahuapan (4.48). En el extremo bajo aparecen Nanacamilpa (2.50), Calpulalpan
(3.00) y Singuilucan (3.08). Sin embargo, Apan, Tepeapulco y Nanacamilpa tienen
sólo una etiqueta observada; sus medias no son estimaciones estables. Municipio
puede ser útil como variable y como grupo diagnóstico, pero una codificación
directa puede sobreajustar categorías pequeñas.

\subsection{Área y rendimiento}
No se observa que las parcelas grandes tengan sistemáticamente mayor rendimiento
por hectárea. La correlación de Spearman es aproximadamente $-0.25$, débil a
moderada y negativa. Parte de esa asociación puede deberse a que el tamaño de las
parcelas también cambia entre estados. El área debe conservarse como predictor,
pero su efecto debe evaluarse controlando ubicación y sin interpretarlo como
causal.

\section{Implicaciones para el modelo}
\begin{enumerate}
  \item Incluir estado, municipio, coordenadas del centroide, elevación y pendiente
  como candidatos, evitando depender únicamente de categorías administrativas.
  \item Comparar validación aleatoria con validación agrupada por municipio o por
  bloques espaciales. Una partición aleatoria puede sobreestimar la capacidad de
  generalizar a otra zona.
  \item Reportar MAE por estado, municipio y grupo de cobertura Landsat, además
  del error global.
  \item No imputar el rendimiento de las 59 parcelas durante entrenamiento. Son
  el conjunto final de inferencia, no valores faltantes comunes.
  \item Mantener t/ha como escala inicial y usar abril--octubre de 2025 como ciclo
  objetivo, cuidando que las variables respeten la fecha de corte de predicción.
\end{enumerate}

\end{document}
